\begin{center}
[2.]\par\medskip
\textsc{Gauss} an \textsc{Bessel}.\quad Göttingen, 26.~Dezember 1821.
\end{center}

\par\noindent\hspace*{\parindent}.\ .\ .\ .\ .\ .\quad Den vergangenen Sommer habe ich größtenteils mit den Vorbereitungen zur Triangulation und der Triangulation selbst verbracht.\quad.\ .\ .\ .\quad Die Winkel an der Sternwarte, am Meridianzeichen, auf dem Hohehagen, dem Hils und dem Brocken sind gemessen. Nur am letzten Punkt ist die Messung des Winkels zwischen dem \textls[120]{Inselsberg} und anderen Punkten missglückt, weil das unerhört schlechte Wetter, solange der Heliotrop dort war, fast alle Beobachtungen unmöglich machte.
Nur einmal konnte die Richtung für eine halbe Stunde und mit dürftigem Ergebnis gemessen werden. Das Ergebnis stimmt jedoch mit anderen früheren, ebenso dürftigen Messungen nicht gut überein, bei denen einmal das Haus und das andere Mal \textsc{Enckes} Sextanten-Vizeheliotrop 3 Minuten vor Sonnenuntergang anvisiert wurde.

\par\noindent\hspace*{\parindent}Dass ich nur einen Heliotrop zur Verfügung hatte, hielt die Arbeit ungemein auf.
Während meines Aufenthalts auf dem Hils musste der Heliotrop nacheinander von Lichtenberg zum Meridianzeichen, zum Brelingerberg und zum Deister gebracht werden. Ebenso wanderte er, als ich auf dem Brocken war, von Lichtenberg zum Hils, zum Hohehagen und zum Inselsberg.
Die wenigsten Winkel konnten also unmittelbar gemessen werden. Ich hatte auf diesen beiden
\par\clearpage
\noindent350\hfill{\small\textbf{BRIEFWECHSEL MIT BESSEL.}}\par\medskip
\noindent Stationen mehrere andere, an sich sichtbare Punkte ausgewählt und maß immer, was sich gerade messen ließ.
Aus dem 4.~Band der Base du système métrique sehe ich, dass die Franzosen es in Spanien ebenso gemacht haben. Sie haben die Messungen aber nicht richtig miteinander verbunden.
Wenn ich im nächsten Jahr meine Triangulation fortsetze, wird es besser gehen, da ich zwei weitere Heliotrope haben werde.
Meinen zum Vizeheliotrop eingerichteten Sextanten habe ich immer bei mir getragen und für telegrafische Zeichen verwendet.
Das funktioniert ganz vortrefflich. Da jener Vizeheliotrop aber auf Winkelabstände unter etwa \(138^{0}\) beschränkt ist, konnte ich ihn nicht immer zu jeder Tageszeit verwenden. Im nächsten Jahr würde ich zu diesem Zweck in der Regel einen der Heliotrope bei mir behalten.

\par\noindent\hspace*{\parindent}Die größte Entfernung, auf die das Heliotroplicht mit bloßem Auge gesehen wurde, ist die vom Brocken zum Hohehagen. Dort wurden meine auf dem Brocken gegebenen telegrafischen Zeichen auf diese Weise gesehen. Die Entfernung beträgt \(9\frac{1}{3}\) geografische Meilen, doch sind dafür wohl günstige Umstände erforderlich.
Bei Entfernungen bis zu 6 Meilen habe ich das Heliotroplicht mit bloßem Auge ohne Mühe gesehen, wenn die Umstände nur einigermaßen günstig waren (trotz des Lichtverlusts beim Gebrauch der Lorgnette). Ebenso sah ich in den Vormittagsstunden das Heliotroplicht vom Hils zum Brocken (\(7\frac{1}{2}\) Meilen).
Die Spiegelfläche beträgt \(2\frac{1}{2}\) Quadratzoll, bei den beiden neuen Heliotropen habe ich \(6\frac{1}{4}\) Quadratzoll gewählt. Das Licht lässt sich leicht abschwächen, doch kann bei großen Entfernungen ein stärkeres Licht manchmal willkommen sein. So drang eines Tages das Licht vom Brelingerberg zum Hils nur kurze Zeit durch, \textls[120]{obwohl} (oder, genauer gesagt, \textls[120]{weil}) der Himmel völlig wolkenlos war.
An solchen Tagen ist die Durchsichtigkeit der Atmosphäre in der Regel am geringsten. Ich erinnere mich an einen solchen Tag, an dem das Heliotroplicht vom Deister zwar recht gut zu sehen war, aber nicht verwendet werden konnte, weil keiner der übrigen Gegenstände zu sehen war, nicht einmal das nur wenig über eine Meile entfernte Einbeck.

\par\noindent\hspace*{\parindent}Das Heliotroplicht hat einen sehr wesentlichen Vorzug vor jedem anderen Signal. Dazu habe ich oft bemerkenswerte Beobachtungen gemacht. Je stärker man vergrößert, desto besser sieht man das Heliotroplicht, irdische Signale hingegen (bei großen Entfernungen) desto schlechter. Denn bei diesen erschwert vor allem ihre \textls[120]{Blässe} das Sehen. Von Hannover habe ich
\par\clearpage
\noindent{\small\textbf{ZUR HANNOVERSCHEN TRIANGULATION.}}\hfill351\par\medskip
\noindent z.~B.\ auf dem Brocken nie eine Spur sehen können, obwohl ich die genaue Lage nach Höhe und Azimut kannte.
An dem Tag, an dem ich das schwache Heliotroplicht vom Inselsberg empfing (\textls[120]{bloß} eine Folge des Wetters, denn bei einigermaßen günstiger Luft hätte das Heliotroplicht bei \(2\frac{1}{2}\) Quadratzoll Spiegelfläche noch mehr als hell genug sein müssen), sah ich es zwar im Fernrohr des Theodoliten noch recht gut, konnte damit aber keine Spur vom Umriss des Berges erkennen. Mein Gehilfe hingegen sah diesen Umriss in einem an sich sehr schlechten Fernrohr mit schwacher Vergrößerung einigermaßen gut, konnte damit aber nur selten das Heliotroplicht sehen.
Das lässt sich leicht erklären. Selbst bei stärkster Vergrößerung bleibt das Heliotroplicht ein Punkt, und sein Licht bleibt immer gleich, doch der Hintergrund ist umso dunkler, je stärker man vergrößert usw.

\par\noindent\hspace*{\parindent}Große Signale habe ich nur zwei gebaut, auf dem Hohehagen und dem Hils. Vermutlich wäre auch das unterblieben, wenn mein Heliotrop sechs Wochen früher fertig gewesen wäre.
Vom Brocken aus zeichnet sich das Hils-Signal vor einem nahen dunklen Hintergrund ab. Während des ganzen Monats, den ich auf dem Brocken verbrachte, habe ich es überhaupt nur zweimal \textls[120]{sehen} können und nur etwa zwei Minuten lang so gut, dass ich einen Winkel hätte messen können. Das Hohehagen-Signal zeichnet sich vor den entfernteren Casseler Bergen ab, war öfter zu sehen und auch ein paar Mal zu beobachten, obwohl ich auch meine Winkelmessung nicht vollständig hätte durchführen können, wenn ich den Heliotrop nicht dorthin geschickt hätte.
Abgesehen von der schlechten Sichtbarkeit, den Kosten und dem Zeitaufwand haben die Signaltürme noch eine sehr unglückliche Seite. Sie bieten einen Anreiz, sie aus rohem Mutwillen zu zerstören.
Leider ist mein Hohehagen-Signal seit Kurzem fast ganz verwüstet, und ich werde froh sein, wenn ich den Punkt auch nur mit hinreichender Genauigkeit wiederfinden kann.

